\scanpage{76}
% archive sheet number: 75
\deleted{satisfaites pour} \inserted{$\unsure{v|\upsilon|u}\in A$ tel que} $\Delta v^{\deleted{\unsure{k}}2}=\deleted{\unsure{y^2-4c}}$, \deleted{\unsure{satisfasse} \illegible{2}} \deleted{(Cela signifie} \inserted{Ceci implique} que
\[
\Delta_0^{\deleted{\unsure{k}}}v_0^2=\deleted{\unsure{\Delta_0^k b_0^2}}\,{b'_0}^{2}
\]
on aura, en supposant que \textcircled{\scriptsize b} \fbox{$\Delta_0$ est él. régulier de $A_0$} que
\[
\left(\frac{b'_0}{\deleted{\unsure{\Delta_0^k}}v_0}\right)^2=\Delta_0\qquad
\text{relation dans }{A_0}_{\Delta_0}
\]
supposons que \fbox{$A_0$ soit intégralemt clos}\textcircled{\scriptsize c} alors on aura $b'_0/\deleted{\unsure{\Delta_0^k}}v_0\in A_0$, soit $b$, donc
\[
b'_0=\deleted{\unsure{\Delta^k}}b_0v_0\qquad(b_0\in A_0,\ b_0^2=\Delta_0).
\]
\deleted{(ce qui ne fait \unsure{aucun sens}}

Bien entendu, on a alors, si $b\equiv b_0\ (2)$ donc $b'=\deleted{\unsure{\Delta^k}}bv$ congru à $b'_0\ (2)$
\[
\gamma(b'_0)\equiv\underset{\substack{\shortparallel\\\deleted{\unsure{\Delta^{2k}}} b^2v^2}}{ {b'}^2}\ \deleted{\unsure{=}}(\bmod 4)
\]
et il reste à exprimer
\[
\deleted{\unsure{\Delta^{2k}}}b^2v^2\equiv\Delta^{\deleted{\unsure{2k}}+1}v^2\quad(4)
\]
on aura, si $\deleted{\unsure{\Delta v|\Delta_0}}\inserted{v\bmod4}$ est régulier dans $A/4A$. -- ce qui est le cas si \fbox{$2$ est régulier dans $A$}\textcircled{\scriptsize d} et \fbox{$v_0$ est régulier dans $A_0$} \deleted{compte tenu que $\Delta_0$ est régulier dans $A_0$}
\[
b^2\deleted{\unsure{v^2}}\equiv\Delta\quad(4)
\]
\deleted{donc}
